\scanpage{136}
\mspage{PU 51}

\margin{\handwritten{\unsure{$\times$|*}}}\inserted{\handwritten{à}} quelles fins...), parmi d'autres structures, d'une soi-disante structure de graphe (dite ``structure d'affinité''), sans qu'il soit trop sûr à aucun moment si telle paire de sommets représente bien ``une arête'' (i.e. si ses deux termes sont considérés ``voisins''), ou non !

Mais une telle situation ne semblera pas tellement étrange au mathématicien \margin{\handwritten{*}}(disons) qu\deleted{\unsure{e}}\handwritten{i} serait rompu à la tâche d'édifier des théories, où les n o t i o n s même avec lesquelles on sera amené à travailler restent encore dans les limbes du non-cré\leavevmode\rlap{e}\handwritten{\'{\phantom{e}}}. Il s'agit\unsure{,}alors de les tailler sur pièces justement, une à une et patiemment, pour arriver à donner un sens à quelque informe nuage d'intuitions, lesquelles peuvent toutes sembler évanescentes et impalpables, mais dont on s e n t pourtant, avec une ``évidence'' pour ainsi dire charnelle et au delà de tout doute, une texture tangible et une chaude substance.

C'est le non-né alors qui nous souffle à chaque moment, au fur et à mesure que le travail de la naissance progresse, quelle doit être telle portion de forme qui s'apprête à naître, et par quel bout la prendre pour la voir émerger du néant et ê t r e. Le tâtonnement même de la main qui saisit la chose encore non née pour l'amener au jour, n'est indécision ni errance, mais prise de connaissance d'où toute hésitation, toute perplexité est absente.

Et ces choses que nous écrivons comme si u n\quad a u t r e écrivait par notre main, et que nous croyons apprendre en les écrivant - quelque part en nous, en des profondeurs ignorées, elles étaient s u e s déjà bien avant que notre main ne les écrive, et n'attendaient que l'attention intense du s c r i b e à l'écoute, qui veuille bien les consigner.

16. La Fleur et son mouvement - ou : plus je m'éloigne, plus je m'approche

(25 mars) Quelques commentaires encore sur notre fleur cosmique, avant de la quitter pour poursuivre le fil interrompu de la réflexion.

Les douze termes yang, placés sur le cercle extérieur, forment aussi les point\deleted{s}\handwritten{es} des douze pétales de la corolle : \deleted{\unsure{c|C}}\handwritten{\unsure{c|C}}eux-ci se rejoignent deux par deux, en les douze points d'insertion sur le cercle intérieur, figurant les douze termes yin de l' ``harmonium''.

Chaque pétale, regardé séparément, se présente aussi comme une sorte de ``mont'' de forme ogivale, dont le sommet yang forme couple avec chacun des deux termes yin, représentés par les points les plus bas\inserted{\handwritten{(*)}} de l'un et l'autre versant du mont. Ces points marquent en même temps le fond des ``vallées'', ou ``ravins'',

\handwritten{(*) Il est entendu ici que la direction vers le ``bas'' est celle vers le centre de la Fleur cosmique.}
% archive sheet number: 136
